\documentclass[10pt,a4paper]{amsart}
\usepackage[margin=19mm]{geometry}
\usepackage{amsmath,amssymb,mathtools}
\usepackage[scr=rsfs,cal=euler]{mathalfa}
\usepackage{xfrac}
\usepackage[unicode,hidelinks,bookmarks=false]{hyperref}
\newcommand{\ie}{i.e.\ }
\newcommand{\cf}{cf.\ }
\newcommand{\resp}{respectively\ }
\newcommand{\Subsection}{\subsection}
\providecommand{\nobreakdash}{\nobreak\hbox{-}}
\setlength{\emergencystretch}{2em}
\multlinegap0pt
\usepackage{amssymb}
\usepackage{mathtools}
\usepackage{braket}
\usepackage{xcolor}
\usepackage{float}
\usepackage{bm}
\newtheorem{lemma}{Lemma}
\title{Operator Algebras and Third Quantization}
\author{Yidong Chen, Marius Junge, Nima Lashkari}
\date{Source: arXiv:2509.02293v2 (CC BY 4.0)}
\begin{document}
\maketitle

\noindent\emph{Source and licence.} Operator Algebras and Third Quantization. Version arXiv:2509.02293v2 (CC BY 4.0), distributed by arXiv under the Creative Commons Attribution 4.0 International licence, http://creativecommons.org/licenses/by/4.0/, as recorded in the arXiv licence metadata for this exact version and shown on the abstract page https://arxiv.org/abs/2509.02293v2. Full licence notice: \texttt{licenses/arxiv-2509.02293-CC-BY-4.0.txt}.\par\smallskip

\noindent\textit{Editorial scope.} Counted unit: the article's lemma lemma (source0.tex lines 2370--2375). The article's complete original proof of it is delivered with it (source0.tex lines 2376--2406, one \texttt{proof} environment of 31 lines). No mathematics is written, repaired or re-derived: the statement and the proof are the article's own lines, in the article's own notation, and no retained line is substituted or re-flowed.  No further source-local context is needed: every cross-reference of every retained line resolves inside the two delivered spans.  Nothing was omitted from any retained span, so the package carries no omission record.  No retained line contains a citation, so the wrapper carries no bibliography and no bibliography entry.  No retained line contains a figure environment, an image-inclusion command or drawing code, so no image is delivered and the package declares no asset dependency.  Editorial formatting only, in addition: an AMS \texttt{amsart} preamble that restates the article's own declarations and macros used by the retained lines, namely the environments \texttt{lemma} on separate counters, together with the standard packages those lines need.  The retained environments print in this document as Lemma 1 in order of appearance, whereas the article prints the same items with the numbering of its own full document.  The complete native source of the article as distributed in the arXiv e-print for this exact version is bundled unchanged as \texttt{sources/184311/source0.tex}.
\begin{lemma}
    \begin{equation}
        \sum_{\sigma\in\mathcal{P}_n}\prod_{A\in\sigma}\sum_{\pi_A\in \mathcal{S}_{|A|}}\mu = \sum_{\pi\in\mathcal{S}_n}B_{|\pi|}(\mu)
    \end{equation}
    where $\mathcal{P}_n$ is the set of partitions of $n$-elements, $\sigma$ is a partition, and $\mathcal{S}_{|A|}$ is the permutation group on $|A|$ elements.
\end{lemma}

\begin{proof}
    For any partition $\sigma =\sqcup A\in\mathcal{P}_n$, we have a set map:
    \begin{equation}
        f_\sigma:\sqcup_A\mathcal{S}_{|A|}\rightarrow \mathcal{S}_n: (\pi_A)_A\mapsto \prod_A\pi_A
    \end{equation}
    Consider the combined set map:
    \begin{equation}
        f = \sqcup_\sigma f_\sigma: \sqcup_\sigma(\sqcup_A\mathcal{S}_{|A|})\rightarrow \mathcal{S}_n
    \end{equation}
    This map is many-to-one. For each fixed permutation $\pi \in\mathcal{S}_n$, consider its cycle decomposition. Let $|\pi|$ be the number of cycles in $\pi$. The cycle decomposition defines a set map:
    \begin{equation}
        \pi:[n]\rightarrow[|\pi|]
    \end{equation}
    where $\pi(i) = \pi(j)$ if and only if $i,j$ are in the same cycle. Then any partition $\eta$ of $[m]$ uniquely induces a partition $\pi^{-1}(\eta)$ of $[n]$ via pullback of $\pi$. In addition, for each subset $A$ in the induced partition $\pi^{-1}(\eta)$ of $[n]$, the cycles of the permutation $\pi$ uniquely determines a permutation $\pi_A\in\mathcal{S}_{|A|}$. Moreover, the product:
    \begin{equation}
        \prod_{A\in \pi^{-1}(\eta)}\pi_A = \pi
    \end{equation}
    Therefore, the preimage under the map $f$ is given by:
    \begin{equation}
        f^{-1}(\pi) = \{(\pi_A)_{A\in \pi^{-1}(\eta)}:\eta\in\mathcal{P}_{|\pi|}\text{ , and }\pi_A\text{ is product of cycles in }\pi\}
    \end{equation}
    We emphasize that once the partition $\eta\in\mathcal{P}_{|\pi|}$ is fixed, the permutation in the preimage $f^{-1}(\pi)$ is uniquely defined by the cycle decomposition of $\pi$.
    Then we have:
    \begin{equation}\label{equation:resummation}
        \sum_{\sigma\in\mathcal{P}_n}\prod_{A\in\sigma}\sum_{\pi_A\in\mathcal{S}_{|A|}}\mu = \sum_{\sigma\in\mathcal{P}_n}\sum_{(\pi_A)_A\in\sqcup_A\mathcal{S}_{|A|}}\mu^{|\sigma|} = \sum_{\pi\in\mathcal{S}_n}\sum_{\eta\in \mathcal{P}_{|\pi|}}\mu^{|\eta|} = \sum_{\pi\in\mathcal{S}_n}B_{|\pi|}(\mu)
    \end{equation}
    In the last equation, we have used the following formula of Touchard polynomial:
    \begin{equation}
        B_n(\mu) = \sum_{\sigma\in\mathcal{P}_n}\mu^{|\sigma|}\qedhere
    \end{equation}
\end{proof}

\end{document}
